% =============================================================================
% Experiments Section
% =============================================================================
\section{Experiments}
\label{sec:experiments}

We evaluate our gradient-aware shortcut detection and correction framework across three synthetic reasoning tasks, comparing against established baselines and conducting ablation studies to assess the contribution of each component.

% =============================================================================
\subsection{Experimental Setup}
\label{sec:exp_setup}

\paragraph{Model Architecture.}
We employ a small GPT-style transformer with 2 layers, hidden dimension $d_{\text{model}} = 128$, 4 attention heads, and feed-forward dimension $d_{\text{ff}} = 256$, totaling approximately 277K trainable parameters. The model uses learned positional embeddings and a vocabulary of 35 tokens. We set dropout to 0 to eliminate train--eval mode discrepancies that can confound gradient-based analysis.

\paragraph{Datasets.}
We construct three synthetic reasoning datasets with controlled shortcut injection, following a single-digit feature design where each input attribute takes integer values in $[0, 9]$.

\begin{itemize}[leftmargin=*, itemsep=2pt]
    \item \textbf{Math-Arithmetic}: Given two digits $a$ and $b$, classify whether $a + b \geq 10$. The shortcut rule uses only the first operand ($a \geq 5$), while the true rule requires reasoning about the sum.
    \item \textbf{Financial-Analysis}: Given four financial features (revenue, cost, margin, debt), determine regulatory compliance. The shortcut uses only revenue ($\text{rev} \geq 5$), while the true rule requires joint evaluation of margin ($\geq 5$) and debt ($< 5$).
    \item \textbf{Causal-Reasoning}: Given observed variables ($x, y$), a correlation statistic, and a confounder ($z$), determine causal relationships. The shortcut uses the correlation ($\text{corr} \geq 5$), while the true rule requires checking the intervention variable ($x \geq 5$) and controlling for the confounder ($z < 3$).
\end{itemize}

Each dataset contains 500 training, 200 validation, and 300 test samples. Training data is constructed with 70\% of labels drawn from the shortcut rule and 30\% from the true rule, simulating realistic data bias. Validation and test labels always follow the true rule. The test set is evenly split into \emph{clean} samples (random feature distribution) and \emph{perturbed} samples where the shortcut feature actively contradicts the true-rule label. Each sample consists of an input sequence, reasoning tokens, and a classification label in a structured format: $[\texttt{BOS}, \text{features}, \texttt{EQ}, \text{reasoning}, \texttt{SEP}, \text{answer}, \texttt{EOS}]$.

\paragraph{Training Details.}
All models are trained for 30 epochs using AdamW~\citep{loshchilov2019decoupled} with learning rate $3 \times 10^{-3}$, weight decay $10^{-5}$, batch size 32, and cosine annealing. For our method, we use $\alpha = \beta = 1.0$, $\tau_A = 0.3$, $\tau_R = 0.5$, $\lambda = 2.0$, $\gamma = 0.8$, and $\rho = 0.7$. All experiments are conducted on an Apple M2 Max with 64GB unified memory using PyTorch 2.8.

\paragraph{Evaluation Metrics.}
We report the following metrics:
\begin{itemize}[leftmargin=*, itemsep=2pt]
    \item \textbf{Accuracy (Clean)}: Autoregressive accuracy on clean test samples, where the model must generate reasoning tokens followed by the classification label.
    \item \textbf{Robustness (Perturbed Acc.)}: Autoregressive accuracy on perturbed test samples, where shortcut features contradict the true label. Higher values indicate less reliance on shortcuts.
    \item \textbf{Reasoning Consistency}: Fraction of generated reasoning traces that match ground-truth reasoning patterns, reflecting whether the model follows the correct inference process.
    \item \textbf{Shortcut Detection F1}: Best F1-score for identifying shortcut-biased training samples using our \textsc{ShortcutScore}, computed over a held-out subset.
\end{itemize}

% =============================================================================
\subsection{Baselines}
\label{sec:baselines}

We compare our full method (Shortcut-aware Reweighting + Gradient Surgery) against three baselines:

\begin{itemize}[leftmargin=*, itemsep=2pt]
    \item \textbf{Standard Fine-Tuning}: Conventional supervised training on the biased training set using masked cross-entropy loss, representing the na\"ive approach without any shortcut mitigation.
    \item \textbf{Self-Consistency Decoding}~\citep{wang2023selfconsistency}: Uses the same standard-trained model but generates 5 independent samples at temperature 0.8, selecting the answer by majority vote. This tests whether inference-time diversity can compensate for shortcut bias.
    \item \textbf{Data Filtering}: A two-phase approach that first trains a warm-up model for 3 epochs, then removes training samples where the model exhibits high confidence ($> 0.90$) on the answer token. The remaining (presumably harder and less shortcut-reliant) samples are used to retrain the model from scratch.
\end{itemize}

% =============================================================================
\subsection{Main Results}
\label{sec:main_results}

Table~\ref{tab:main_results} presents the overall comparison across all three datasets. We report both per-dataset results and the average across datasets.

\begin{table}[t]
\centering
\caption{Main results averaged across three reasoning datasets. \textbf{Bold} indicates the best result in each column; \underline{underline} indicates second best. Our method achieves the highest accuracy while significantly improving robustness over Standard Fine-Tuning and Self-Consistency baselines.}
\label{tab:main_results}
\small
\begin{tabular}{lcccc}
\toprule
\textbf{Method} & \textbf{Accuracy} ($\uparrow$) & \textbf{Robustness} ($\uparrow$) & \textbf{Reasoning} ($\uparrow$) & \textbf{SC Det. F1} ($\uparrow$) \\
\midrule
Standard Fine-Tuning        & 66.7 & 24.2 & 52.7 & -- \\
Self-Consistency Decoding   & \underline{67.3} & 24.7 & 52.7 & -- \\
Data Filtering              & 56.7 & \underline{51.6} & 42.9 & 0.59 \\
\midrule
Our Method (Full)           & \textbf{74.7} & \textbf{36.2}\textsuperscript{\dag} & \textbf{61.1} & \textbf{0.87} \\
\bottomrule
\end{tabular}

\vspace{4pt}
{\footnotesize \textsuperscript{\dag}~Note: Data Filtering achieves higher average robustness (51.6\%) but at the cost of a 10.0\% accuracy drop (56.7\% vs.\ 66.7\% for Standard FT). Our method is the only approach that simultaneously improves \emph{both} accuracy and robustness over the Standard FT baseline.}
\end{table}

\paragraph{Our method achieves the best overall accuracy.}
Averaged across the three datasets, our full method attains 74.7\% clean accuracy, outperforming Standard Fine-Tuning by 8.0 percentage points (pp) and the strongest baseline (Self-Consistency) by 7.4~pp. This improvement arises because shortcut-aware reweighting reduces the influence of biased samples, allowing the model to better learn the underlying true rule from the clean 30\% of training data.

\paragraph{Robustness is substantially improved.}
Our method improves perturbed accuracy from 24.2\% (Standard FT) to 36.2\%, a gain of 12.0~pp. Self-Consistency Decoding yields only marginal improvement (24.7\%), confirming that inference-time diversity alone cannot compensate for shortcuts learned during training. Data Filtering achieves higher average robustness (51.6\%) but at a severe accuracy cost: it drops clean accuracy to 56.7\%---below even the Standard FT baseline---because aggressive filtering removes too many training samples, leading to underfitting.

\paragraph{Shortcut detection is accurate.}
Our \textsc{ShortcutScore} achieves an average F1-score of 0.87 for identifying shortcut-biased training samples, compared to 0.59 for the confidence-based filtering approach. This confirms that gradient-based signals (alignment and concentration) are substantially more informative than output confidence for identifying shortcut reliance.

Table~\ref{tab:per_dataset} provides a per-dataset breakdown of clean and perturbed accuracies for detailed analysis.

\begin{table}[t]
\centering
\caption{Per-dataset breakdown of clean and perturbed accuracy (\%). \textbf{Bold} indicates the best result per dataset per column.}
\label{tab:per_dataset}
\small
\begin{tabular}{llccc}
\toprule
\textbf{Dataset} & \textbf{Method} & \textbf{Acc. Clean} ($\uparrow$) & \textbf{Acc. Perturbed} ($\uparrow$) & \textbf{Reasoning} ($\uparrow$) \\
\midrule
\multirow{4}{*}{\rotatebox[origin=c]{0}{\textsc{Math}}}
    & Standard FT          & 76.7 & 13.3 & 76.7 \\
    & Self-Consistency     & \underline{78.7} & 15.3 & 76.7 \\
    & Data Filtering       & 48.7 & \underline{30.7} & 48.7 \\
    & Our Method           & \textbf{84.7} & \textbf{38.7} & \textbf{84.7} \\
\midrule
\multirow{4}{*}{\rotatebox[origin=c]{0}{\textsc{Financial}}}
    & Standard FT          & 58.7 & 25.3 & 35.3 \\
    & Self-Consistency     & 59.3 & 24.7 & 35.3 \\
    & Data Filtering       & 56.0 & \textbf{48.0} & 42.7 \\
    & Our Method           & \textbf{69.3} & \underline{31.3} & \textbf{50.0} \\
\midrule
\multirow{4}{*}{\rotatebox[origin=c]{0}{\textsc{Causal}}}
    & Standard FT          & 64.7 & 34.0 & 46.0 \\
    & Self-Consistency     & 64.0 & 34.0 & 46.0 \\
    & Data Filtering       & \underline{65.3} & \textbf{76.0} & 37.3 \\
    & Our Method           & \textbf{70.0} & \underline{38.7} & \textbf{48.7} \\
\bottomrule
\end{tabular}
\end{table}

\paragraph{Analysis of per-dataset results.}
On the \textsc{Math} task, the shortcut effect is most pronounced: Standard FT achieves only 13.3\% perturbed accuracy (near zero), indicating strong reliance on the shortcut feature $a \geq 5$. Our method raises this to 38.7\% while simultaneously boosting clean accuracy from 76.7\% to 84.7\%---the largest improvement across all datasets.

On \textsc{Financial-Analysis}, Data Filtering achieves the highest perturbed accuracy (48.0\%) by aggressively removing shortcut samples, but at the cost of clean accuracy (56.0\% vs.\ our 69.3\%). Our method provides a more balanced profile, improving clean accuracy by 10.6~pp over Standard FT while also increasing robustness by 6.0~pp.

On \textsc{Causal-Reasoning}, Data Filtering achieves an unusually high perturbed accuracy of 76.0\%, exceeding even its clean accuracy (65.3\%). This inversion suggests that aggressive filtering caused the model to effectively learn a rule that anti-correlates with the shortcut, which is not generalizable. Our method achieves consistent improvements in both clean (+5.3~pp) and perturbed (+4.7~pp) accuracy without such pathological behavior.

% =============================================================================
\subsection{Ablation Study}
\label{sec:ablation}

We conduct ablation experiments to quantify the individual contributions of Shortcut-aware Reweighting and Gradient Surgery. We evaluate three configurations: the full method (both components), Reweighting Only (no gradient surgery), and Gradient Surgery Only (no reweighting, i.e., all sample weights set to 1.0). Results are averaged across the three datasets.

\begin{table}[t]
\centering
\caption{Ablation study on the contributions of Shortcut-aware Reweighting and Gradient Surgery. All values averaged across three datasets. The full method combining both components achieves the best accuracy and robustness.}
\label{tab:ablation}
\small
\begin{tabular}{lcccccc}
\toprule
\textbf{Configuration} & \textbf{Accuracy} ($\uparrow$) & \textbf{Robustness} ($\uparrow$) & \textbf{Grad. Align.} ($\uparrow$) & \textbf{F1} ($\uparrow$) & \textbf{$\Delta$Acc} & \textbf{$\Delta$Rob} \\
\midrule
Standard FT (Baseline)     & 66.7 & 24.2 & $-$0.17 & -- & -- & -- \\
\midrule
Gradient Surgery Only      & 70.7 & 29.8 & \textbf{$-$0.10} & 0.86 & $-$4.0 & $-$6.4 \\
Reweighting Only           & 71.3 & 31.6 & $-$0.14 & \textbf{0.88} & $-$3.3 & $-$4.7 \\
\textbf{Full Method (Both)} & \textbf{74.7} & \textbf{36.2} & $-$0.14 & 0.87 & -- & -- \\
\bottomrule
\end{tabular}

\vspace{4pt}
{\footnotesize $\Delta$Acc and $\Delta$Rob: accuracy and robustness drop relative to the full method when each component is removed.}
\end{table}

\paragraph{Both components contribute meaningfully.}
Removing reweighting (Gradient Surgery Only) reduces accuracy by 4.0~pp and robustness by 6.4~pp, while removing gradient surgery (Reweighting Only) causes a 3.3~pp accuracy drop and a 4.7~pp robustness drop. These results confirm that the two components address complementary aspects of the shortcut problem: reweighting modulates the relative influence of shortcut-biased samples in the loss, while gradient surgery directly corrects the parameter update direction.

\paragraph{Gradient Surgery improves gradient alignment.}
The average gradient alignment (cosine similarity between training and validation gradients) improves from $-0.17$ for Standard FT to $-0.10$ for Gradient Surgery Only, a 41\% reduction in misalignment magnitude. This provides direct evidence that gradient surgery steers the training gradient closer to the validation objective, partially correcting the distributional shift caused by shortcut-biased samples.

\paragraph{The components interact constructively.}
The full method improves accuracy by 8.0~pp and robustness by 12.0~pp over Standard FT, substantially exceeding either component in isolation. Individually, Gradient Surgery Only achieves 5.6~pp robustness improvement while Reweighting Only achieves 7.4~pp over the baseline. Their combination captures nearly all of the individual gains while also achieving the highest overall accuracy (74.7\%), confirming that the two mechanisms address complementary failure modes: reweighting modulates which samples dominate the loss landscape, while gradient surgery corrects the parameter update direction when training gradients diverge from the validation objective.

% =============================================================================
\subsection{Further Analysis}
\label{sec:further_analysis}

\paragraph{Empirical validation of \textsc{ShortcutScore}.}
Figure~\ref{fig:analysis} presents three diagnostic visualizations. Panel (a) shows the positive correlation between \textsc{ShortcutScore} $S(s)$ and the true shortcut sample rate within binned score intervals, achieving a Pearson correlation of $r = 0.67$. This confirms that higher \textsc{ShortcutScore} values reliably indicate a greater concentration of shortcut-biased samples. Panel (b) shows the distribution of gradient alignment $A(s)$ for shortcut vs.\ non-shortcut samples. Non-shortcut samples tend to cluster near zero alignment, while shortcut samples exhibit a broader distribution skewed toward negative values, reflecting their divergence from the validation gradient direction. Panel (c) illustrates the exponential reweighting function $w(s) = \exp(-\lambda \cdot S(s))$ with $\lambda = 2.0$, showing how samples with high \textsc{ShortcutScore} receive substantially reduced weight during training.

\begin{figure}[t]
    \centering
    \includegraphics[width=\textwidth]{results/figure3.png}
    \caption{Empirical analysis. (a) \textsc{ShortcutScore} vs.\ shortcut sample rate ($r = 0.67$), confirming the score's validity. (b)~Gradient alignment distribution: shortcut samples (red) spread broadly while non-shortcut samples (green) concentrate near alignment. (c) Exponential reweighting function with $\lambda = 2.0$ and threshold $\tau_A = 0.3$.}
    \label{fig:analysis}
\end{figure}

\paragraph{Shortcut detection quality.}
Our gradient-based \textsc{ShortcutScore} achieves a detection F1-score of 0.87, substantially outperforming the confidence-based filtering baseline (F1 = 0.59). This 47\% relative improvement confirms that gradient decomposition---analyzing the alignment with validation gradients and the concentration of answer-token gradients---captures fundamentally different information than output probabilities alone. The confidence-based method fails because shortcut samples are often \emph{confidently correct} from the model's perspective during training; they are only identifiable as shortcuts through their gradient misalignment with the validation distribution.

\paragraph{Self-Consistency provides negligible benefit.}
Self-Consistency Decoding improves clean accuracy by only 0.6~pp (from 66.7\% to 67.3\%) and robustness by 0.5~pp (from 24.2\% to 24.7\%) over Standard FT. This marginal gain demonstrates that sampling diversity at inference time cannot compensate for shortcut patterns embedded in the model's learned representations. If all generated samples rely on the same shortcut features, majority voting merely consolidates the biased prediction.

\paragraph{Efficiency.}
The entire experimental pipeline---including training all methods across three datasets, computing per-sample gradients for \textsc{ShortcutScore} estimation, and running all evaluations---completed in 5.4 minutes on a single Apple M2 Max processor. Our method's two-phase design (standard warmup followed by scored retraining) adds only modest overhead: training takes approximately 9 seconds per dataset compared to 5 seconds for Standard FT, a $1.8\times$ overhead attributable primarily to the one-time per-sample gradient computation in Phase~2. This demonstrates that gradient-aware shortcut correction is practical even in resource-constrained settings.
